% !TeX root = ../main.tex
% Section 8: Conclusion
% Batch #4.8 — Post M1+M2 axis, notation hygiene
% v4.0 revision (2026-08-30): F1-F6 narrative-closure micro-edits applied.
%   F1 two-pillar demotion of Lindblad; F2 15-day claim demoted to additional
%   signal; F3 VIX first-differences backstop; F4 polynomial-time restatement;
%   F5 risk-regime-barometer terminology restored; F6 Weyl-Heisenberg dash
%   unified with Introduction. Abstract and Introduction remain axis-frozen.

\section{Conclusion}\label{sec:conclusion}

This paper develops a non-commutative quantum probability framework for option
pricing under market microstructural friction. The framework rests on two
pillars: the Weyl-Heisenberg algebra
$[\hat{X}, \hat{P}] = i\hbar_{\mathrm{econ}} \mathbf{1}$ and the GNS Hilbert
space pricing formula $P = \langle H^* | \Phi \rangle e^{-rT}$; the market
state's temporal evolution is governed by a Lindblad master equation
(Section~\ref{sec:weyl-heisenberg}). The economic Planck constant
$\hbar_{\mathrm{econ}} = \lambda \cdot q_0$ links the framework to Kyle's price
impact parameter, providing a direct microstructural interpretation.

The theoretical contributions are threefold. First, variance-optimal hedging
under price impact is equivalent to a quantum measurement, with the optimal
hedge given by the orthogonal projection of the payoff onto the hedgeable
subspace (Theorem~\ref{thm:existence-hedge}, Theorem~\ref{thm:hedging-measurement}).
Second, the leading-order quantum correction to the implied volatility surface,
derived via WKB expansion (Theorem~\ref{thm:effective-vol}), produces a universal
$\mathcal{O}(\tau^{-1})$ short-maturity skew divergence---distinct from the
bounded skew of Black--Scholes--Heston and the $\mathcal{O}(\tau^{H-1/2})$
scaling of rough volatility. Third, the non-commutative pricing formula
interpolates between the classical Black--Scholes price in the commutative
limit $\hbar_{\mathrm{econ}} \to 0$ and a quantum-corrected price with
$\mathcal{O}(\hbar_{\mathrm{econ}})$ corrections in the presence of
microstructural friction.

The empirical results yield three findings. First, the quantum coherence
$\mathcal{C}(\rho_t)$, defined as the $\ell_1$-norm of the off-diagonal
elements of the reconstructed density matrix, provides a forward-looking
early-warning signal that leads the classical Pearson correlation between
KRE and VIX by several weeks for gradually developing failures and by a
contemporaneous response for sudden ones---as documented in the 2023 US
regional banking crisis case study (Silicon Valley Bank on March~10,
Signature Bank on March~12, and First Republic Bank on May~1)---and by a
comparable margin over the VIX first-differences benchmark
(Section~\ref{sec:additional-results}). Second, as an additional
entropy-based signal, the von Neumann entropy
$S(\rho_t)$ crosses its 95th-percentile baseline threshold 15~trading days
(approximately three weeks) before the SVB failure of March~10, 2023,
providing an even earlier warning signal that Granger-causes changes in
KRE option-implied volatility. Third, the quantum optimal hedge from
Proposition~\ref{prop:quantum-optimal-hedge} reduces hedging errors relative
to both Black--Scholes and Heston benchmarks during a stress-window case study, with the
improvement concentrated in out-of-the-money regimes where the microstructural
correction $\hbar_{\mathrm{econ}} \Gamma_{\text{BS}} \gamma x / \tau$ is
non-negligible.

The empirical implementation of these results relies on Quantum State
Tomography (QST): the daily density operator $\rho_t$ is reconstructed
from observable macro-financial moments via a strictly convex semidefinite
program that is solvable in polynomial time via off-the-shelf convex
optimization solvers, yielding a positive-semidefinite, unit-trace state
without heuristic parameter calibration.

Several directions merit further work. Extending the single-qubit
parametrization to multi-qubit or continuous-variable representations is
conceptually natural but raises computational challenges. Developing dynamic,
data-driven mappings from macro-finance observables to Pauli coefficients
would enhance adaptability across asset classes. Out-of-sample validation on
the 2008 Global Financial Crisis (with an alternative credit-spread
observable, given HY-OAS data limitations), the 2011 European Debt Crisis,
and the 2020 COVID-19 pandemic would test generalizability. Extending the
framework to American options, exotic derivatives, and multi-asset structured
products represents a further direction for both theory and practice.

The non-commutative quantum probability framework addresses three structural
puzzles in derivatives pricing: operator ordering ambiguity under price
impact, the $\mathcal{O}(\tau^{-1})$ short-maturity skew divergence, and the
inability of conventional risk frameworks to capture systemic contagion. In
practical use, the framework operates as a forward-looking risk-regime
barometer, complementing rather than replacing existing volatility- and
correlation-based indicators. The results suggest that non-commutative methods
offer a promising direction for bridging mathematical finance theory with the
empirical complexity of financial markets under microstructural constraints.
